\chapter{Evaluation}
\label{ch:evaluation}

Chapter~\ref{ch:methodology} fixed everything that produces a retrieved list or a generated
summary in this thesis: the four retrieval systems, the shared MMR mechanism they all select
through, the two use cases and the operating points each is evaluated at, the hardened-pool
construction that stresses retrieval under adversarial distractor pressure, and the dual-annotator
protocol that turns every candidate argument into a graded relevance judgment. What remains is to
turn a retrieved list, or a summary generated from one, into the numbers Chapter~\ref{ch:results}
reports and compares across systems. That is the subject of this chapter.

The two use cases give rise to two largely disjoint metric families, because they ask a retrieval
system to produce two different kinds of output. UC1's retrieved lists are read directly, so its
metrics ask the questions a moderator scanning the list would ask: are the most relevant arguments
near the top, is the set free of near-duplicates, does it contain genuine pro/con tension rather
than merely topic-adjacent material, and, once the pool is deliberately hardened, how much of what
came back is contamination by design. UC2's output is a generated summary rather than a list, so
its metrics ask a different family of questions: does the summary say what the retrieved arguments
actually said, does it invent anything the retrieved arguments do not support, and does it
represent both sides of the debate fairly. Section~\ref{sec:evaluation-uc1} defines the first
family and Section~\ref{sec:evaluation-uc2} the second.

Within each family, one subset of metrics forms the primary comparison this thesis reports across
systems, matching the metric set reported in the conference paper this thesis expands on. A second
subset is retained as supplementary diagnostics: metrics computed throughout evaluation-protocol
development that surfaced a property of the metrics themselves worth documenting in its own right
(most notably in Section~\ref{sec:evaluation-uc1-diversity}), or that this thesis develops in more
depth than the original paper did. Both subsets are defined here in full, with the primary or
supplementary status of each metric stated explicitly rather than left for the reader to infer from
which table it eventually appears in.

Before either family can be computed, every retrieved argument has to be resolved to a
ground-truth grade. Section~\ref{sec:evaluation-matching} defines that resolution step first, since
nearly every metric that follows in Section~\ref{sec:evaluation-uc1} depends on it.

\section{Resolving Retrieved Arguments to Ground-Truth Grades}
\label{sec:evaluation-matching}

Every UC1 relevance metric, and UC2's Coverage metric, needs a way to look up a retrieved
argument's grade in the dual-annotator ground truth of
Section~\ref{sec:methodology-groundtruth}. For a system whose candidate pool is drawn directly
from a policy's own graded arguments, this lookup is close to trivial: the retrieved text is one
of the exact strings the annotators graded, so an exact, whitespace-stripped string match against
the ground truth resolves the large majority of retrievals for StandardRAG, RAG-no-MMR, and
Random. PR-GraphRAG is the case that requires more care. Its one-hop \textsc{contradicts}
expansion (Section~\ref{sec:methodology-prgraphrag}) can, under the hardened pool, pull in a
distractor argument that was graded as part of the distractor manifest rather than the target
policy's own pool (Section~\ref{sec:methodology-groundtruth-distractors}), and, more rarely under
the native pool, a genuine same-policy argument that the 200-item truncation happened to exclude
from what was graded. An exact match against the target policy's own graded pool then fails, not
because the argument is unjudged, but because it was judged under a different key.

A retrieved argument that fails exact match is resolved by falling back to semantic matching: it
is embedded with the same \textsc{sentence-transformers/all-mpnet-base-v2} model used throughout
this thesis \cite{reimers2019sbert,song2020mpnet}, compared by cosine similarity against every
argument the retrieving system had graded access to, and assigned the grade of its single closest
match if that similarity clears $\tau_{\text{match}} = 0.75$, or treated as unmatched (grade 0, no
credit) otherwise. An earlier version of this evaluation pipeline skipped the fallback entirely and
resolved every retrieval by a plain substring lookup, so any retrieved argument phrased as even a
close paraphrase of a graded one, of which the corpus has many, silently scored as a complete miss
regardless of how relevant it actually was. This inflated the apparent gap between systems for
reasons that had nothing to do with retrieval quality, since some systems happened to retrieve more
verbatim pool members than others rather than because their retrieved arguments were less relevant,
and it was diagnosed and corrected before any of the results in Chapter~\ref{ch:results} were
produced. Two design choices in the resulting fallback matter for what the metrics defined later in
this chapter actually measure.

\begin{itemize}
\item \textbf{Exact match is tried first, semantic match only as a fallback.} Every argument in the
judged pool is looked up verbatim before any embedding is computed, so resolution for the three
systems that never leave the judged pool is exact and introduces no matching noise at all. The
embedding-based fallback exists specifically for the minority of retrievals that can fall outside
the pool a policy's own arguments were graded from, and is applied only to those, rather than
replacing exact lookup everywhere and paying its similarity-threshold uncertainty on every
retrieval regardless of whether it is needed.

\item \textbf{Resolution is single-best-match, not nearest-any-relevant-item.} A retrieved argument
is credited with the grade of the one ground-truth argument it is closest to, not the highest grade
found among any ground-truth argument within the similarity threshold. This closes what would
otherwise be a saturation loophole on self-similar pools: without single-best-match resolution, a
retrieved argument sitting near a cluster containing one highly-graded item could inherit that
item's grade even if the retrieved argument itself is a weaker, less specific restatement of the
same point, inflating every relevance metric that depends on the resolved grade.
\end{itemize}

Every relevance-dependent metric defined in Section~\ref{sec:evaluation-uc1} is written in terms of
the grade this resolution step assigns, written $g(d)$ for a retrieved argument $d$, together with
a match type, exact, semantic, or unmatched, that this thesis reports as a diagnostic alongside the
relevance metrics themselves: a system whose top-$k$ list resolves mostly by exact match is
retrieving verbatim pool members, while a rising semantic-match fraction signals that a growing
share of its retrievals are reaching outside the graded pool, which is expected of PR-GraphRAG
under the hardened pool and is otherwise a signal worth inspecting on its own.

\section{UC1 Metrics: Conflict-Map Retrieval}
\label{sec:evaluation-uc1}

Table~\ref{tab:uc1-metrics} summarizes the eleven quantities computed over every UC1 retrieved
list. Five of them, nDCG@$k$, MeanRel@$k$, Conflict density, EILD, and Distractor-rate@$k$, are
the primary metrics this thesis reports as its headline UC1 comparison across systems, matching
the metric set reported in the underlying conference paper. The remaining six, Precision@$k$,
Recall@$k$, S-Recall@$k$, raw ILD, the LLM rubric judge, and the abandoned pairwise judge, are
supplementary: computed throughout evaluation-protocol development, retained in the full evaluation
report, and used in this chapter and in Chapter~\ref{ch:results} either as interpretive diagnostics
or, in the case of raw ILD, specifically to expose a property of the primary metric it was replaced
by. This design continues a tradition of graded, rank-aware relevance evaluation for argument
retrieval established by shared evaluation tasks in the area \cite{touche2022,perspectivearg2024},
adapted here to a setting where every candidate pool, not just a system's top-$k$ output, has been
exhaustively graded (Section~\ref{sec:methodology-groundtruth-protocol}).

\begin{table}[htbp]
\centering
\caption{UC1 metrics at a glance. Status marks whether a metric is part of this thesis's primary
reported comparison or a supplementary diagnostic. Pool marks whether a metric is meaningful under
the native pool, the hardened pool, or both.}
\label{tab:uc1-metrics}
\begin{tabular}{llll}
\toprule
\textbf{Metric} & \textbf{Status} & \textbf{Range} & \textbf{Pool} \\
\midrule
nDCG@$k$ & Primary & $[0,1]$ & Native \& hardened \\
MeanRel@$k$ & Primary & $[0,3]$ & Native \& hardened \\
Conflict density & Primary & $[0,1]$ & Native \& hardened \\
EILD & Primary & $[0,1]$ & Native \& hardened \\
Distractor-rate@$k$ & Primary & $[0,1]$ & Hardened only \\
Precision@$k$ & Supplementary & $[0,1]$ & Native \& hardened \\
Recall@$k$ & Supplementary & $[0,1]$ & Native \& hardened \\
S-Recall@$k$ & Supplementary & $[0,1]$ & Native \& hardened \\
ILD (raw) & Supplementary & $[0,1]$ & Native \& hardened \\
LLM rubric judge (3 axes) & Supplementary & $[1,5]$ & Native \& hardened \\
Pairwise judge & Abandoned & -- & -- \\
\bottomrule
\end{tabular}
\end{table}

\subsection{Primary Relevance Metrics: nDCG@$k$ and MeanRel@$k$}
\label{sec:evaluation-uc1-relevance}

The primary relevance metric is graded nDCG@$k$ \cite{jarvelin2002ndcg}, computed with exponential
rather than linear gain,

\begin{equation}
\mathrm{DCG}@k = \sum_{i=1}^{k} \frac{2^{g(d_i)} - 1}{\log_2(i+1)}, \qquad
\mathrm{nDCG}@k = \frac{\mathrm{DCG}@k}{\mathrm{IDCG}@k},
\label{eq:evaluation-ndcg}
\end{equation}

\noindent where $g(d_i) \in \{0,1,2,3\}$ is the resolved grade of the argument at rank $i$
(Section~\ref{sec:evaluation-matching}) and $\mathrm{IDCG}@k$ is the DCG of the ideal ranking, the
$k$ highest consensus grades available in the judged pool, sorted descending. Exponential gain
rewards placing a grade-3 argument above a grade-2 one far more heavily than linear gain would,
which matches how UC1 is actually read: a moderator scanning a five-item conflict map benefits far
more from having the single most essential argument at the top than from a list that is merely,
on average, somewhat more relevant. Because the native and hardened pools are both graded
exhaustively rather than only over a positional shortlist
(Section~\ref{sec:methodology-groundtruth-protocol}), $\mathrm{IDCG}@k$ is well defined for every
policy and stance, an argument that would belong in the ideal top-$k$ can never be missing from the
denominator simply because it fell past wherever a shortlist boundary happened to be drawn.

nDCG@$k$ is reported at both operating points fixed in Chapter~\ref{ch:methodology}, $k=5$ under
the native pool and $k=5$ under the hardened pool, and Chapter~\ref{ch:results} additionally
reports it across the pool-size and $k$ sweep introduced there to check whether conclusions drawn
at the single $k=5$ operating point generalize. A linear-gain variant, $\mathrm{DCG}@k = \sum_i
g(d_i) / \log_2(i+1)$, is retained as a robustness check rather than reported as a second primary
metric: it exists to confirm that any ordering observed under exponential gain is not an artifact
of the specific gain function chosen, not because it measures something the exponential variant
does not already capture.

MeanRel@$k$ is nDCG@$k$'s rank-free companion,

\begin{equation}
\mathrm{MeanRel}@k = \frac{1}{k} \sum_{i=1}^{k} g(d_i),
\label{eq:evaluation-meanrel}
\end{equation}

\noindent the mean resolved grade of the top-$k$ on the original 0--3 scale, dividing always by
$k$ rather than by the number of items actually returned, so a system cannot raise its score simply
by returning fewer arguments than requested. MeanRel@$k$ sits between a binary relevance metric,
which collapses the full ordinal grade scale to a single relevant/non-relevant cut, and nDCG@$k$,
which uses the full grade scale but folds it together with rank position. Reporting it alongside
nDCG@$k$ answers a question nDCG@$k$ alone cannot: whether an observed difference between two
systems reflects a genuine difference in the grades their retrieved arguments carry, or only a
difference in how those grades happen to be ordered within the list. Two systems that retrieve the
identical multiset of grades in different orders receive different nDCG@$k$ scores but the same
MeanRel@$k$, so a gap that appears in nDCG@$k$ but closes in MeanRel@$k$ is a ranking difference,
while a gap that survives in both is a difference in what was retrieved, not merely how it was
arranged.

\subsection{Rank-Free Diagnostics: Precision@$k$, Recall@$k$, and Subtopic Recall}
\label{sec:evaluation-uc1-diagnostics}

Three further relevance-based quantities are computed as supplementary diagnostics, each answering
a narrower question than nDCG@$k$ or MeanRel@$k$ and each retained for a specific interpretive
reason rather than as redundant coverage of the same construct.

Precision@$k$ uses the binary relevant/non-relevant cut already fixed in
Section~\ref{sec:methodology-groundtruth-consensus}, grade $\geq 2$,

\begin{equation}
\mathrm{Precision}@k = \frac{1}{k} \left| \{\, i \leq k : g(d_i) \geq 2 \,\} \right|,
\label{eq:evaluation-precision}
\end{equation}

\noindent resolved, like every relevance metric in this chapter, by single-best-match
(Section~\ref{sec:evaluation-matching}): a retrieved item counts as a hit only if its own resolved
grade clears the threshold, not because it happens to sit near some other relevant item in
embedding space. Precision@$k$ is easier to read at a glance than a graded metric, and, because it
depends only on the binary cut rather than the full grade scale, it is the metric most directly
comparable to the Distractor-rate@$k$ metric of Section~\ref{sec:evaluation-uc1-distractorrate},
both being simple fractions of the top-$k$ that satisfy a yes/no condition.

Recall@$k$ is the fraction of the judged pool's relevant arguments (grade $\geq 2$, deduplicated by
ground-truth text) that appear anywhere in the top-$k$,

\begin{equation}
\mathrm{Recall}@k = \frac{\left| \{\, a \in \text{pool} : g(a) \geq 2 \,\} \cap \{d_1,\dots,d_k\} \right|}
{\left| \{\, a \in \text{pool} : g(a) \geq 2 \,\} \right|}.
\label{eq:evaluation-recall}
\end{equation}

\noindent Full-pool judging, rather than judging only a shortlist or only the union of what
compared systems retrieve, is what makes this denominator well defined at all, following the same
incomplete-judgment concern documented for pooled retrieval evaluation more generally
\cite{buckley2004bpref}: a fixed, judged relevant set that every system is measured against, rather
than one collected only from what was retrieved, avoids penalizing a system for surfacing a
legitimate relevant argument that a narrower judging protocol never saw. Recall@$k$ is retained for
continuity with earlier evaluation passes, but this thesis reports it as a demoted diagnostic
rather than a primary metric, because its denominator is structurally uninformative at the
operating points UC1 is actually reported at. The native pool's relevance base rate is 57.7\%
(Section~\ref{sec:methodology-groundtruth-agreement}), so a typical 200-item pool contains on the
order of 115 relevant arguments per stance, meaning even a perfect top-5 retrieval can never clear
roughly $5/115 \approx 0.04$ Recall@5, and the hardened pool's lower 38.4\% base rate
(Section~\ref{sec:methodology-groundtruth-distractors}) still leaves a denominator in the same
range. Recall@$k$ at $k=5$ or $k=15$ is therefore compressed into a narrow band near zero for every
system regardless of retrieval quality, and any small differences that do appear are dominated by
noise in which few relevant items happen to fall in the top-$k$ rather than by a meaningful
difference in coverage.

S-Recall@$k$ is this thesis's response to that problem, following the subtopic-retrieval framing of
\textcite{zhai2003srecall}: rather than asking what fraction of individually relevant arguments a
system surfaces, it asks what fraction of the distinct \emph{points} present in the relevant set a
system surfaces at least one instance of. The ground truth stores only a $(\text{text}, \text{grade})$
pair per argument, with no subtopic label, so a notion of subtopic has to be operationalized before
S-Recall@$k$ can be computed at all. This thesis operationalizes it by clustering the judged pool's
relevant arguments, per policy and stance, with agglomerative average-linkage clustering over
cosine distance between \textsc{sentence-transformers/all-mpnet-base-v2} embeddings, at a distance
threshold of 0.15 (cosine similarity $\approx 0.85$) and no fixed number of clusters, since how many
distinct argumentative points a given policy's relevant set actually contains is not known in
advance and should vary from one policy to the next. S-Recall@$k$ is then the fraction of the
resulting clusters represented at least once, by a relevant, resolved retrieval, in the top-$k$:

\begin{equation}
\mathrm{S\text{-}Recall}@k = \frac{\left|\, \{\, c(a) : a \in \{d_1,\dots,d_k\},\ g(a) \geq 2 \,\}\, \right|}
{n_{\text{subtopics}}},
\label{eq:evaluation-srecall}
\end{equation}

\noindent where $c(a)$ denotes the subtopic cluster a relevant retrieved argument $a$ resolves to
and $n_{\text{subtopics}}$ is the total number of clusters found for that policy and stance.
Clustering the relevant set into subtopics this way is a defensible operationalization rather than
the only possible one, and the specific distance threshold is a genuine methodological choice
rather than a value with an independent justification of its own. A sensitivity check across two or
three neighboring thresholds would strengthen the claim that S-Recall@$k$'s conclusions are not an
artifact of exactly where that threshold was set, and is noted here as a limitation of the current
operationalization rather than carried out in this thesis.

\subsection{Diversity: From Intra-List Distance to Relevance-Constrained Diversity}
\label{sec:evaluation-uc1-diversity}

Diversity is not an afterthought in UC1's evaluation design, it is the second half of exactly the
tradeoff MMR is built to manage (Section~\ref{sec:methodology-mmr}), so a metric that measures it
directly is needed alongside the relevance metrics of
Sections~\ref{sec:evaluation-uc1-relevance}--\ref{sec:evaluation-uc1-diagnostics}. Diversity as a
first-class retrieval objective has a long history in recommendation and IR evaluation
\cite{ziegler2005diversification,clarke2008novelty} and continues to be treated as one in recent
retrieval-augmented generation work \cite{lu2026diversity}. The most direct, reference-free way to
measure it is intra-list distance,

\begin{equation}
\mathrm{ILD} = 1 - \operatorname*{mean}_{i<j} \mathrm{sim}(e_i, e_j),
\label{eq:evaluation-ild}
\end{equation}

\noindent one minus the mean pairwise cosine similarity over every pair of items in the retrieved
set, following \textcite{carbonell1998mmr}. ILD is exactly the quantity MMR's redundancy term
(Equation~\ref{eq:mmr}) greedily works against at selection time, so it is the natural first choice
for measuring what that term bought a system, and this thesis reports it as a supplementary
diagnostic for that reason.

Raw ILD, however, breaks in a specific and instructive way once the pool is hardened. Injected
distractor arguments are drawn from topically adjacent sibling policies
(Section~\ref{sec:methodology-hardened-pipeline}), close enough to the target policy to be
plausible near-misses, but they are not close to the target policy's own genuine arguments in
embedding space in the way two genuine arguments on the same specific point tend to be. A system
that retrieves several such distractors alongside a few genuine arguments therefore scatters its
retrieved set widely across embedding space almost by construction, since the distractors sit far
from the genuine cluster and often far from each other as well, and raw ILD rewards exactly this
scattering regardless of whether any of it is relevant. A retrieval system could, in principle,
score near the top of the ILD distribution by retrieving five mutually unrelated, entirely
irrelevant arguments, which is the opposite of what UC1's diversity criterion is meant to reward: a
conflict map that spans several genuinely distinct points of contention, not one that merely spans
several unrelated topics.

EILD, an expected, relevance-constrained variant of intra-list distance, is the primary diversity
metric this thesis reports for exactly this reason. It restricts the pairwise computation of
Equation~\ref{eq:evaluation-ild} to the subset of the retrieved list that resolves to a relevant
grade,

\begin{equation}
\mathrm{EILD} = 1 - \operatorname*{mean}_{\substack{i<j \\ g(d_i) \geq 2,\, g(d_j) \geq 2}}
\mathrm{sim}(e_i, e_j),
\label{eq:evaluation-eild}
\end{equation}

\noindent returning 0 when fewer than two relevant items are retrieved, which occurs occasionally
at $k=5$ and becomes rare at the wider $k=15$ operating point UC2 draws from. A grade-0 distractor
is excluded from every pair EILD sums over, so a retrieval set padded with irrelevant, embedding-
distant filler can no longer inflate its diversity score the way it can inflate raw ILD. What EILD
measures instead is genuinely the construct UC1 cares about: among the arguments a moderator would
actually credit as relevant, how varied are the distinct points they make. This thresholded
construction follows the broader relevance-aware turn in diversity and novelty evaluation
\cite{ziegler2005diversification,clarke2008novelty}, restricting the diversity computation itself to
a relevant subset rather than treating diversity and relevance as two independent criteria, though
it stops short of a fully relevance-weighted formulation in which every pair's contribution would be
scaled continuously by both items' grades rather than gated by a single threshold. This thesis
reports raw ILD alongside EILD, labeled explicitly as a diagnostic, specifically so that the
artifact just described remains visible in the full evaluation report rather than silently
disappearing once EILD replaced it as the headline diversity number.

Concretely, suppose a $k=5$ retrieved set resolves to grades $3, 2, 0, 2, 0$. Raw ILD averages over
all $\binom{5}{2}=10$ pairs, including every pair that touches one of the two grade-0 items, so a
system that happened to retrieve two embedding-distant distractors at ranks 3 and 5 sees its ILD
pulled upward by six of those ten pairs regardless of how similar the three genuinely relevant items
are to one another. EILD averages only over the $\binom{3}{2}=3$ pairs formed by the grade-3 and
two grade-2 items, so if those three genuine arguments happen to restate a similar point, EILD
reports that redundancy directly instead of letting the two irrelevant items mask it.

\subsection{Distractor-Rate@$k$}
\label{sec:evaluation-uc1-distractorrate}

Distractor-rate@$k$ is a direct provenance metric, meaningful only under the hardened pool, that
counts the fraction of the top-$k$ originating from the injected distractor manifest
(Section~\ref{sec:methodology-hardened-pipeline}) rather than from the target policy's own native
argument set,

\begin{equation}
\mathrm{DR}@k = \frac{\left| \{\,d_1,\dots,d_k\,\} \cap \text{distractor manifest} \right|}{k}.
\label{eq:evaluation-distractorrate}
\end{equation}

\noindent Unlike every relevance metric in this chapter, Distractor-rate@$k$ needs no resolved
grade at all: it asks a question about where a retrieved argument came from, not about how relevant
it is, and the distractor manifest recorded at injection time (Section~\ref{sec:methodology-hardened-pipeline})
answers that question directly. This makes Distractor-rate@$k$ an independent cross-check on the
grade-based relevance metrics rather than a restatement of them in different units. A distractor is
not, by construction, guaranteed to be irrelevant to the target policy, a sibling-policy argument
drawn for topical closeness can occasionally be graded relevant on its own merits
(Section~\ref{sec:methodology-groundtruth-distractors}), so a low Distractor-rate@$k$ and a high
nDCG@$k$ are correlated but not identical claims: the first says a system's structural filtering
kept contamination out of the retrieved list, the second says what it did retrieve was well graded,
and a system could in principle score well on one without the other.

\subsection{Conflict Density}
\label{sec:evaluation-uc1-conflictdensity}

Every metric defined so far asks whether a retrieved argument is relevant to its policy. Conflict
density asks a different question entirely, one that goes directly to what UC1 is actually for: does
the retrieved PRO set engage the retrieved CON set, or does it merely sit alongside it. A retrieved
list can be highly relevant by every grade-based metric above while still failing UC1's purpose if
its pro and con arguments talk past each other rather than contesting the same underlying claims,
since a conflict map that does not show genuine conflict has not done its job regardless of how
relevant each individual argument is. Conflict density is the mean natural-language-inference
contradiction probability over every retrieved PRO$\times$CON pair,

\begin{equation}
\mathrm{ConfDensity} = \operatorname*{mean}_{p \in \text{PRO}, c \in \text{CON}} P(\text{contradict} \mid p, c),
\label{eq:evaluation-confdensity}
\end{equation}

\noindent computed by the identical \textsc{cross-encoder/nli-deberta-v3-base} classifier
\cite{he2023debertav3} introduced in Section~\ref{sec:methodology-edges} for \textsc{contradicts}
edge recovery, applied here to retrieved argument pairs rather than to candidate corpus pairs. The
two uses of this NLI classifier differ in what they do with its output. Edge recovery converts a
contradiction probability into a binary accept/reject decision, requiring the bidirectional minimum
of forward and backward scores to clear $\tau_{\text{contra}} = 0.90$
(Section~\ref{sec:methodology-edges-calibration}), because a graph edge is either present or absent.
Conflict density instead reports the continuous mean probability directly, with no threshold and no
bidirectionality requirement, because it is a diagnostic score meant to be averaged and compared
across systems rather than a structural decision that has to be made once and stored. A retrieved
set can therefore register a moderate conflict density even where no individual pair would have
cleared the edge-recovery threshold, which is expected: conflict density is measuring the degree of
argumentative tension present in what was retrieved, not counting how many \textsc{contradicts}
edges happen to already exist between those specific arguments in the graph.

\subsection{LLM Rubric Judging and the Abandoned Pairwise Comparison}
\label{sec:evaluation-uc1-rubric}

The five metrics of Sections~\ref{sec:evaluation-uc1-relevance}--\ref{sec:evaluation-uc1-conflictdensity}
are all decomposed, verifiable quantities: each one reduces to a well-defined computation over
resolved grades, embeddings, or NLI scores, with a narrow and inspectable failure mode. This thesis
also computes a holistic quality judgment as a supplementary cross-check, following the
LLM-as-judge line of evaluation \cite{zheng2023judging,liu2023geval}, on the reasoning that a
decomposed metric suite can miss a failure only a reader of the full retrieved set would notice, and
that agreement between the holistic judgment and the decomposed metrics is itself evidence that the
decomposed metrics are measuring what they claim to.

The first design tried for this holistic judgment was a binary, double-sided pairwise comparison: an
LLM judge shown two systems' retrieved sets for the same policy and asked which one it preferred,
each pair judged in both presentation orders to cancel positional bias, a precaution motivated by
documented evidence that LLM judges are not immune to such artifacts
\cite{wang2024fair,panickssery2024selfpref}. This design was tried and abandoned. Because the
systems compared in this thesis retrieve from overlapping, often near-identical candidate pools
under the native condition, a forced binary choice produced tie rates in the 73--97\% range across
policies, leaving too little discriminative signal to be useful for anything beyond confirming that
the judge could not reliably tell two similar retrieved sets apart.

The rubric judge that replaced it scores each system's retrieved set independently rather than
comparatively, on a 1--5 scale across three named axes, using Llama 3.3 70B \cite{grattafiori2024llama3}
at temperature 0 for reproducibility:

\begin{itemize}
\item \textbf{Conflict engagement.} Do the PRO and CON arguments directly engage each other,
attacking the same premises and dimensions, rather than talking past each other, the same construct
Conflict density (Section~\ref{sec:evaluation-uc1-conflictdensity}) targets through NLI rather than
holistic judgment.

\item \textbf{Redundancy avoidance.} Does each argument in the set add a distinct point, with
little repetition within it, the same construct EILD (Section~\ref{sec:evaluation-uc1-diversity})
targets through embedding distance rather than holistic judgment.

\item \textbf{Balanced writing aid.} Would this retrieved set equip a citizen to write a balanced,
well-informed position paper covering the strongest pros, the strongest cons, and the underlying
value tensions between them, a construct none of the decomposed metrics above measures directly.
\end{itemize}

Scoring each set independently rather than pairwise restores meaningful spread, since ties then
occur only where two sets are judged genuinely equal on an axis rather than being forced whenever a
comparative judgment is too close to call. The rubric judge is reported as a supplementary
cross-check rather than promoted into the primary metric set, both because it lacks the narrow,
verifiable failure mode of the decomposed metrics above and because its usefulness depends on
whether it actually separates systems, a question Section~\ref{sec:evaluation-uc2-rubric} returns
to directly for UC2's version of the same judge, where the answer turns out to be less favorable
than it is here.

\section{UC2 Metrics: Debate Summarization}
\label{sec:evaluation-uc2}

UC1's metrics all operate on a retrieved list. UC2 asks a different question of the same
underlying retrieval: once a system's top-$k$ arguments are handed to a summarizer, does the
resulting summary preserve what made that retrieval good or bad, or does the generation step wash
those differences out. This question is not answerable from UC1's metrics alone, since retrieval
quality and downstream generation quality are related but distinct constructs that need not move
together \cite{samuel2026beyond}, which is exactly why this thesis evaluates both use cases rather
than treating UC1 as a sufficient proxy for UC2. Summaries are generated with Qwen3:8b
\cite{yang2025qwen3} at temperature 0, from the structured five-field prompt fixed in
Section~\ref{sec:methodology-uc2}, and this section defines the four automatic metrics computed
over the resulting summary text.

\subsection{Coverage}
\label{sec:evaluation-uc2-coverage}

Coverage measures whether the summary omits genuine debate content that was available to it,
operationalized as NLI-entailment-based recall of the retrieved argument set,

\begin{equation}
\mathrm{Coverage} = \frac{\left| \{\, a \in \text{retrieved} : \max_{s \in \text{summary}}
P(\text{entail} \mid s, a) > 0.5 \,\} \right|}{|\text{retrieved}|},
\label{eq:evaluation-coverage}
\end{equation}

\noindent the fraction of retrieved arguments for which some sentence in the summary entails it.
Framing coverage as entailment rather than lexical overlap follows the NLI-based line of
summarization evaluation \cite{laban2022summac}, and specifically the recall-oriented framing used
in general summarization evaluation \cite{fabbri2021summeval} and in retrieval-augmented generation
evaluation toolkits \cite{es2024ragas}, adapted here to the retrieved-argument set rather than to a
reference summary: an argument the summary paraphrases or restates counts as covered even where no
words overlap, while an argument the summary never mentions, regardless of how it is worded,
does not. The entailment probability is computed by the same \textsc{cross-encoder/nli-deberta-v3-base}
classifier used throughout this thesis (Section~\ref{sec:methodology-edges},
Section~\ref{sec:evaluation-uc1-conflictdensity}), applied here to summary-sentence and
source-argument pairs rather than to argument-argument pairs, and the $0.5$ threshold is the same
entailment cutoff Section~\ref{sec:evaluation-uc2-balance} reuses for Stance Balance below.

\subsection{Faithfulness}
\label{sec:evaluation-uc2-faithfulness}

Faithfulness measures the opposite failure mode: does the summary assert anything not grounded in
what was actually retrieved. Rather than a holistic judgment of factual accuracy, Faithfulness is
computed as structured, per-claim verification, following the claim-decomposition approach to
generation evaluation \cite{liu2023geval,es2024ragas}. An LLM judge, Llama 3.3 70B at temperature 0,
first extracts every distinct factual or argumentative claim from the summary text, explicitly
excluding mere framing statements that assert nothing substantive about the debate, and then decides
for each extracted claim whether it is supported, directly stated or clearly entailed, by the
retrieved source arguments the summarizer was given. Faithfulness is the resulting fraction,

\begin{equation}
\mathrm{Faithfulness} = \frac{\#\ \text{supported claims}}{\#\ \text{claims in summary}}.
\label{eq:evaluation-faithfulness}
\end{equation}

\noindent This binary, per-claim outcome is what distinguishes Faithfulness from a holistic rubric
score of the kind Section~\ref{sec:evaluation-uc1-rubric} defined for UC1: each claim receives an
independently checkable supported/unsupported verdict against a specific source passage rather than
a single scalar impression of the summary as a whole, closer in kind to Coverage and Precision@$k$
than to the rubric judge, which is exactly why it is reported among UC2's primary four metrics
rather than alongside the rubric judge Section~\ref{sec:evaluation-uc2-rubric} excludes.

\subsection{Stance Balance}
\label{sec:evaluation-uc2-balance}

UC2 was introduced in Section~\ref{sec:methodology-uc2} as helping a citizen understand both sides
of a debate, not only whichever side a system happened to retrieve more of, so a metric that
measures whether the generated summary actually reflects both sides is needed directly rather than
inferred from Coverage or Faithfulness alone. Stance Balance reuses the same NLI entailment
machinery as Coverage, applied separately to the retrieved PRO arguments and the retrieved CON
arguments: for each stance, it counts how many of that stance's retrieved arguments the summary
entails above the same $0.5$ threshold, and computes the proportion of covered content attributable
to the PRO side,

\begin{equation}
\mathrm{pro\_frac} = \frac{\text{covered}_{\text{pro}}}{\text{covered}_{\text{pro}} + \text{covered}_{\text{con}}},
\qquad
\mathrm{StanceBalance} = 1 - 2\,\left| \mathrm{pro\_frac} - 0.5 \right|.
\label{eq:evaluation-balance}
\end{equation}

\noindent A summary that entails an equal number of PRO and CON arguments scores 1.0, perfect
balance, and a summary that entails arguments from only one side scores 0.0, regardless of how many
arguments from that one side it covers. Stance Balance is deliberately independent of Coverage: a
summary can be perfectly balanced while covering very little of either side, or highly one-sided
while covering a great deal of one side and none of the other, and the two metrics are reported
side by side rather than combined into a single number precisely so that this distinction remains
visible.

\subsection{BERTScore F1 and the Independent Reference Summary}
\label{sec:evaluation-uc2-bertscore}

The three metrics above all score the summary against its own retrieved arguments. BERTScore F1
\cite{zhang2020bertscore} instead scores it against an independently generated reference summary,
giving UC2 one metric that is not itself computed from the same source material the summary was
asked to be faithful to. BERTScore matches each candidate-summary token to its most similar
reference-summary token by contextual embedding cosine similarity, and symmetrically each reference
token to its most similar candidate token, and reports the F1 of the two resulting matching scores.
Because the matching is over embeddings rather than surface tokens, a candidate and reference that
paraphrase the same content in entirely different words still score highly, which is the property
that makes BERTScore suitable here: two summaries of the same debate, written by different models,
are unlikely to share much vocabulary even when they agree on substance.

The reference summary itself has to be constructed carefully to avoid a specific circularity. If the
reference were generated by the same model used as the candidate summarizer, Qwen3:8b, BERTScore
would partly reward stylistic self-similarity, a summary would score well for sounding like
something the same model tends to produce, rather than for capturing the same content a genuinely
different model would independently converge on. This thesis therefore generates the reference
summary with a model held distinct from both the candidate summarizer and the LLM used for
Faithfulness judging in Section~\ref{sec:evaluation-uc2-faithfulness}, following the same
disjoint-model reasoning already used for the dual-annotator ground truth
(Section~\ref{sec:methodology-groundtruth-protocol}) and motivated by the same documented
self-preference risk in LLM-based evaluation \cite{panickssery2024selfpref}. A further, less obvious
design choice is that the reference summary is regenerated separately for each pool condition and
each system's own retrieved argument set, rather than produced once as a single global gold
reference reused everywhere. Given a universal reference, BERTScore would conflate two different
questions, whether a system's retrieval enabled a summary close to one fixed target, and whether the
candidate summarizer faithfully rendered whatever it was actually given, since a system that
retrieves a different set of arguments than the reference was built from is penalized for that
difference regardless of how well it summarized its own retrieval. Generating the reference from the
same retrieved arguments the candidate summary was built from isolates the second question instead:
BERTScore then measures whether the candidate summarizer's rendering of a given input diverges from
a stronger model's rendering of the identical input, a generation-consistency check rather than a
retrieval-quality-via-reference check, and the two are not the same claim.

\subsection{Why UC2 Omits Holistic Rubric Judging}
\label{sec:evaluation-uc2-rubric}

UC1's rubric judge (Section~\ref{sec:evaluation-uc1-rubric}) is retained as a useful supplementary
diagnostic because it shows meaningful spread across systems and axes, evidence that it is
discriminating between retrieved sets of genuinely different quality rather than defaulting to a
uniform response. UC2's version of the same judge, an independent 1--5 score across coverage,
balance, insight, and specificity, scored by Llama 3.3 70B at temperature 0 on the same
direct-scoring design that replaced UC1's abandoned pairwise comparison
(Section~\ref{sec:evaluation-uc1-rubric}), was computed during evaluation-protocol development and
did not show the same spread. Every system's mean score on every axis clustered inside a narrow band
near the top of the 1--5 scale, typically between 4.0 and 5.0, regardless of which retrieval system
had produced the summary being judged. A metric that assigns nearly every summary a score within one
point of the ceiling carries almost no information for ranking systems against each other: whatever
genuine quality differences exist between summaries produced from different retrievals, the rubric
judge is not resolving them.

This thesis therefore follows the underlying conference paper's decision and reports UC2's rubric
judge as a supplementary, non-primary diagnostic, in favor of the four automatic metrics of
Sections~\ref{sec:evaluation-uc2-coverage}--\ref{sec:evaluation-uc2-bertscore}, each of which has a
narrower, more verifiable failure mode than a single holistic scalar and, unlike the rubric judge,
demonstrably separates systems in practice. This is a deliberate asymmetry with UC1 rather than an
inconsistency: the same judging methodology behaves differently on the two use cases because UC1's
retrieved sets are heterogeneous enough for a holistic judge to distinguish, while UC2's
five-field structured summaries, produced from a fixed prompt template regardless of which system
fed them, apparently converge in surface quality closely enough that a coarse 1--5 scale cannot tell
them apart even where the automatic metrics still can.

\section{Reporting Conventions and Metric Reliability}
\label{sec:evaluation-reporting}

Every metric defined in this chapter is computed at the level of a single (policy, stance) pair for
UC1 or a single policy for UC2, then aggregated in two stages before it reaches a comparison table
in Chapter~\ref{ch:results}. For UC1's stance-aware metrics, a PRO score and a CON score are first
computed independently for each policy and averaged into one per-policy figure, so that a policy
with an unusually strong PRO side and a weak CON side, or vice versa, does not silently favor
whichever stance a system happens to handle better. Per-policy figures are then averaged, together
with their standard deviation, across the fixed 30-policy stratified sample of
Section~\ref{sec:methodology-policy-selection}. The per-policy figures themselves are retained in
the full evaluation report rather than discarded once the aggregate is computed, both because
Chapter~\ref{ch:results} cross-references them against policy-level covariates such as argument-count
density and \textsc{equivalent}-community modularity, and because retaining a per-unit breakdown
rather than only a pooled mean is standard practice in comparing IR systems across a fixed set of
topics, where paired, per-topic comparison is generally more informative than comparing aggregate
averages alone \cite{smucker2007significance,demsar2006statistical}.

Two further reporting conventions follow directly from decisions already fixed in
Chapter~\ref{ch:methodology}. Every metric is reported separately under the native pool and the
hardened pool, since Section~\ref{sec:evaluation-uc1-diversity} and
Section~\ref{sec:evaluation-uc1-distractorrate} have already shown that a metric's behavior, and in
EILD's case its very validity as a diversity measure, can differ qualitatively between the two
conditions. And every UC1 metric is reported at the fixed $k=5$ operating point as the primary
comparison, with the pool-size and $k$ sweep of Chapter~\ref{ch:results} reported separately as a
generalization check rather than folded into the same table, since a conclusion that only holds at
one specific operating point is a materially weaker claim than one that holds across a swept range.

A metric's reported value is only as trustworthy as its stability across independent evaluation
runs, a question that matters most for the Random baseline, since a random retriever's score on any
given evaluation run is, by construction, one draw from a distribution rather than a deterministic
property of the system. As a validation step taken during evaluation-protocol development, Random's
retrieval was re-run over 20 independent random seeds at the native-pool, $k=5$ operating point, and
nDCG@5 was recomputed on each. The resulting distribution has a mean of 0.388 and a 95\% bootstrap
confidence interval \cite{efron1993bootstrap} of $[0.377, 0.399]$, a width of roughly 0.022 around a
value close to the single-seed figure reported as Random's headline score in
Chapter~\ref{ch:results}. An interval this narrow, and this far from zero, indicates that Random's
elevated apparent competitiveness under the native pool (Section~\ref{sec:methodology-groundtruth-distractors})
is a stable, structural property of the native pool's high relevance density, not an artifact of one
unusually favorable random draw. This distinction matters directly for how Chapter~\ref{ch:results}
interprets any comparison involving Random: a tie between Random and a designed retrieval system
under the native pool, given this check, is evidence that the metric's headroom is compressed on
that pool, an interpretation this thesis returns to when discussing why the pool was hardened in the
first place, rather than evidence that either system's evaluation run happened to be lucky. This kind
of seed-level robustness check is one instance of a broader concern in constructing retrieval
benchmarks that are difficult enough, and stable enough, to actually discriminate between systems
rather than being dominated by evaluation noise \cite{thakur2025freshstack}.